% from ega3-1-fr.tex
\oldpage[III]{101}
Comme $X$ est recouvert par les ouverts affines $D_+(T_i)$ en nombre
$r+1$, (ii) résulte de (1.4.12). Notons d'autre part que (iv) a déjà été
démontré (II, 2.7.9).

Nous allons prouver simultanément (i) et (iii). Notons que ces assertions
sont vraies pour $\mathcal F=\mathcal O_X(m)$
(\hyperref[III.2.1.13-fr]{2.1.13})~; elles le sont donc aussi lorsque
$\mathcal F$ est somme directe d'un nombre fini de
$\mathcal O_X$-Modules de la forme $\mathcal O_X(m_j)$. D'autre part,
(i) et (iii) sont vraies trivialement pour $q>r$ en vertu de (ii). Nous
allons procéder par \emph{récurrence descendante} sur $q$. On sait que
$\mathcal F$ est isomorphe à un quotient d'une somme directe $\mathcal E$
d'un nombre fini de faisceaux $\mathcal O_X(m_j)$ (II, 2.7.10)~;
autrement dit, on a une suite exacte
\[
0\longrightarrow\mathcal R\longrightarrow\mathcal E
\longrightarrow\mathcal F\longrightarrow0,
\]
où $\mathcal R$ est cohérent (0\textsubscript{I}, 5.3.3) et où
$\mathcal E$ vérifie (i) et (iii). Comme $\mathcal F(n)$ est un foncteur
exact en $\mathcal F$, on a aussi la suite exacte
\[
0\longrightarrow\mathcal R(n)\longrightarrow\mathcal E(n)
\longrightarrow\mathcal F(n)\longrightarrow0
\]
pour tout $n\in\mathbb Z$. On en déduit la suite exacte de cohomologie
\[
H^{q-1}(X,\mathcal E(n))\longrightarrow
H^{q-1}(X,\mathcal F(n))\longrightarrow
H^q(X,\mathcal R(n)).
\]
Comme $\mathcal E(n)$ est somme directe des $\mathcal O_X(n+m_j)$
(II, 2.5.14), $H^{q-1}(X,\mathcal E(n))$ est de type fini, et il en est
de même de $H^q(X,\mathcal R(n))$ par l'hypothèse de récurrence~; comme
$A$ est noethérien, on en conclut que
$H^{q-1}(X,\mathcal F(n))$ est de type fini pour tout $n\in\mathbb Z$, et
en particulier pour $n=0$. D'autre part, par l'hypothèse de récurrence,
il existe un entier $N$ tel que pour $n\geq N$ on ait
$H^q(X,\mathcal R(n))=0$~; par ailleurs, on peut supposer aussi $N$ choisi
tel que $H^{q-1}(X,\mathcal E(n))=0$ pour $n\geq N$, puisque $\mathcal E$
vérifie (iii)~; on en conclut que
$H^{q-1}(X,\mathcal F(n))=0$ pour $n\geq N$, ce qui achève la
démonstration.
\end{proof}

\begin{corollary}[2.2.3]
\label{III.2.2.3-fr}
Sous les hypothèses de (\hyperref[III.2.2.1-fr]{2.2.1}), soit
$\mathcal F\to\mathcal G\to\mathcal H$ une suite exacte de
$\mathcal O_X$-Modules cohérents. Il existe alors un entier $N$ tel que
pour $n\geq N$, la suite
\[
f_*(\mathcal F(n))\longrightarrow f_*(\mathcal G(n))
\longrightarrow f_*(\mathcal H(n))
\]
soit exacte.
\end{corollary}

\begin{proof}
Soient $\mathcal F'$, $\mathcal G'$, $\mathcal G''$ le noyau, l'image et
le conoyau de $\mathcal F\to\mathcal G$~; $\mathcal G'$ est le noyau et
$\mathcal G''$ l'image de $\mathcal G\to\mathcal H$, soit $\mathcal H''$
le conoyau de cet homomorphisme~; tous ces $\mathcal O_X$-Modules sont
cohérents (0\textsubscript{I}, 5.3.4). Comme $\mathcal F(n)$ est un
foncteur exact en $\mathcal F$, il suffit de prouver que pour $n$ assez
grand, chacune des suites
\begin{align*}
0&\longrightarrow f_*(\mathcal F'(n))
\longrightarrow f_*(\mathcal F(n))
\longrightarrow f_*(\mathcal G'(n))\longrightarrow0,\\
0&\longrightarrow f_*(\mathcal G'(n))
\longrightarrow f_*(\mathcal G(n))
\longrightarrow f_*(\mathcal G''(n))\longrightarrow0,\\
0&\longrightarrow f_*(\mathcal G''(n))
\longrightarrow f_*(\mathcal H(n))
\longrightarrow f_*(\mathcal H''(n))\longrightarrow0
\end{align*}
est exacte~; par suite, on peut supposer que
$0\to\mathcal F\to\mathcal G\to\mathcal H\to0$ est exacte. On a alors
la suite exacte de cohomologie
\[
0\longrightarrow f_*(\mathcal F(n))
\longrightarrow f_*(\mathcal G(n))
\longrightarrow f_*(\mathcal H(n))
\longrightarrow R^1f_*(\mathcal F(n))\longrightarrow\cdots
\]
et la conclusion résulte de (\hyperref[III.2.2.1-fr]{2.2.1}, (ii)).
\end{proof}

\begin{corollary}[2.2.4]
\label{III.2.2.4-fr}
Soient $Y$ un préschéma noethérien, $f:X\to Y$ un morphisme de type fini,
$\mathcal L$ un $\mathcal O_X$-Module inversible ample pour $f$~; pour
tout $\mathcal O_X$-Module $\mathcal F$, on pose
$\mathcal F(n)=\mathcal F\otimes_{\mathcal O_X}\mathcal L^{\otimes n}$
(pour $n\in\mathbb Z$). Soit
$\mathcal F\to\mathcal G\to\mathcal H$ une suite exacte de
$\mathcal O_X$-Modules cohérents, telle que les supports de $\mathcal F$
et de $\mathcal H$ soient propres sur $Y$ (II, 5.4.10). Il existe alors
un entier $N$ tel que, pour $n\geq N$, la suite
\[
f_*(\mathcal F(n))\longrightarrow f_*(\mathcal G(n))
\longrightarrow f_*(\mathcal H(n))
\]
soit exacte.
\end{corollary}

\begin{proof}
Le même raisonnement qu'au début de
(\hyperref[III.2.2.1-fr]{2.2.1}) montre que si le corollaire est vrai
pour $\mathcal L^{\otimes d}$ $(d>0)$, il l'est aussi pour $\mathcal L$~;
on peut donc se borner au cas où $\mathcal L$ est très ample pour $f$
(II, 4.6.11), et par suite on peut identifier $X$ à un ouvert dans un
$Y$-schéma $Z=\operatorname{Proj}(\mathcal S)$, où $\mathcal S$ est une
$\mathcal O_Y$-Algèbre quasi-cohérente graduée à degrés positifs, dans
laquelle $\mathcal S_1$ est de type fini et engendre $\mathcal S$, de
sorte que $\mathcal L=i^*(\mathcal O_Z(1))$, où $i$ est l'immersion
canonique $X\to Z$ (II, 4.4.7). Cela étant, comme
$\operatorname{Supp}(\mathcal G)$ est fermé dans $X$ et contenu dans
$\operatorname{Supp}(\mathcal F)\cap\operatorname{Supp}(\mathcal H)$, il
est propre sur $Y$~; les supports de $\mathcal F$, $\mathcal G$,
$\mathcal H$ sont donc \emph{fermés dans $Z$} (II, 5.4.10). Les faisceaux
$\mathcal F'=i_*(\mathcal F)$, $\mathcal G'=i_*(\mathcal G)$,
$\mathcal H'=i_*(\mathcal H)$ sont donc des $\mathcal O_Z$-Modules
cohérents, et la suite $\mathcal F'\to\mathcal G'\to\mathcal H'$ est
exacte~; en outre, si $g:Z\to Y$ est le morphisme structural, on a
$f=g\circ i$, et il est clair que
$\mathcal F'(n)=i_*(\mathcal F(n))$, et de même pour $\mathcal G'$ et
$\mathcal H'$~; la conclusion résulte donc de
(\hyperref[III.2.2.3-fr]{2.2.3}) appliqué à
$\mathcal F'$, $\mathcal G'$, $\mathcal H'$.
\end{proof}

\begin{env}[2.2.5]
\label{III.2.2.5-fr}
\emph{Remarques.} ---
(i) L'assertion (i) de (\hyperref[III.2.2.1-fr]{2.2.1}) est encore vraie
lorsqu'on suppose seulement que $Y$ est \emph{localement noethérien}~; en
effet, la propriété est évidemment locale sur $Y$~; d'autre part, les
hypothèses de (\hyperref[III.2.2.1-fr]{2.2.1}) impliquent que pour tout
ouvert $U\subset Y$, la restriction de $f$ à $f^{-1}(U)$ est un morphisme
projectif $f^{-1}(U)\to U$ (II, 5.5.5, (iii)) et
$\mathcal L|f^{-1}(U)$ est ample pour ce morphisme (II, 4.6.4).

(ii) L'assertion (iii) de (\hyperref[III.2.2.1-fr]{2.2.1}) est encore
valable, comme on l'a vu, lorsque l'on suppose seulement que $X$ est un
schéma quasi-compact ou un préschéma dont l'espace sous-jacent est
noethérien, et $f:X\to Y$ un morphisme quasi-compact (II, 4.6.8). Mais il
faut noter que même lorsqu'on suppose que $Y$ est le \emph{spectre d'un
corps} $K$ et que $f$ est \emph{quasi-projectif}, l'assertion (ii) de
(\hyperref[III.2.2.1-fr]{2.2.1}) n'est plus nécessairement vérifiée. Par
exemple, soit $X'=\operatorname{Spec}(K[T_0,\ldots,T_r])$ et soit $X$ la
réunion des ouverts affines $D(T_i)$ de $X'$ $(0\leq i\leq r)$~; comme
l'immersion $X\to X'$ est quasi-compacte, le morphisme structural
$f:X\to Y$ est quasi-affine (II, 5.1.10), donc $\mathcal O_X$ est très
ample pour $f$ (II, 5.1.6). Mais l'anneau
$\Gamma(X,\mathcal O_X)$ s'identifie à l'intersection des anneaux de
fractions $(K[T_0,\ldots,T_r])_{T_i}$ pour $0\leq i\leq r$
(I, 8.2.1.1), c'est-à-dire à $K[T_0,\ldots,T_r]$. Par suite, il résulte
des formules (1.4.3.1) et (1.1.3.5) que l'on a
$H^r(X,\mathcal O_X^{\otimes n})=H^r(X,\mathcal O_X)=A\neq0$ pour tout
$n$.
\end{env}

\subsection{Application aux faisceaux gradués d'algèbres et de modules}
\label{subsection:III.2.3-fr}

\begin{theorem}[2.3.1]
\label{III.2.3.1-fr}
Soient $Y$ un préschéma noethérien, $\mathcal S$ une
$\mathcal O_Y$-Algèbre graduée à degrés positifs, quasi-cohérente et de
type fini, $X=\operatorname{Proj}(\mathcal S)$,
$q:X\to Y$ le morphisme structural, $\mathcal M$ un $\mathcal S$-Module
gradué quasi-cohérent vérifiant la condition \textbf{(TF)}. Alors il existe
un entier $N$ tel que, pour $n\geq N$, l'homomorphisme canonique
(II, 8.14.5.1)
\[
\alpha_n:\mathcal M_n\longrightarrow
q_*(\shProj_0(\mathcal M(n)))
=q_*((\shProj(\mathcal M))_n)
\]
